\documentclass[12pt]{article}
\usepackage{amsmath, amssymb, array}

\begin{document}
	
	\section*{The Symmetric Group $S_3$}
	
	\subsection*{1. Definition}
	
	Let 
	\[
	X = \{1,2,3\}.
	\]
	The symmetric group on $X$ is
	\[
	S_3 = Aut_{set}(X) = \mathrm{Sym}(X),
	\]
	the set of all bijections $f: X \to X$.  
	It contains $3! = 6$ elements.
	
	\subsection*{2. Group Operation}
	
	For $f,g \in S_3$, the operation is composition:
	\[
	(g \cdot f)(x) = g(f(x)), \qquad \forall x \in X.
	\]
	
	\subsection*{3. Elements of $S_3$}
	
	We list all six elements, with their action on $(1,2,3)$ and their cycle notation:
	
	\[
	\begin{array}{c|c|c}
		\text{Name} & \text{Mapping on } (1,2,3) & \text{Cycle Notation} \\ \hline
		\text{Identity} & (1,2,3) \mapsto (1,2,3) & () \\ \hline
		f & (1,2,3) \mapsto (2,1,3) & (1\,2) \\ \hline
		f' & (1,2,3) \mapsto (1,3,2) & (2\,3) \\ \hline
		f'' & (1,2,3) \mapsto (3,2,1) & (1\,3) \\ \hline
		g & (1,2,3) \mapsto (3,1,2) & (1\,3\,2) \\ \hline
		g' & (1,2,3) \mapsto (2,3,1) & (1\,2\,3) \\
	\end{array}
	\]
	
	\subsection*{4. Transpositions}
	
	A \emph{transposition} is a cycle of length $2$; it swaps two elements and fixes the third.  
	In $S_3$, the three transpositions are:
	\[
	(1\,2), \quad (2\,3), \quad (1\,3).
	\]
	These generate the entire group $S_3$, since any permutation can be written as a product of transpositions.
	
	\subsection*{Transposition (Formal)}
	
	A \textbf{transposition} on a set \(X\) with \(|X| = n \geq 2\) is the simplest non-trivial permutation: it swaps two elements and leaves all others fixed.  
	
	\medskip
	
	\noindent
	\textbf{Definition.}  
	Let \(f : X \to X\) be a bijection. \(f\) is a transposition if there exist distinct \(x,y \in X\) such that
	\[
	f(x) = y, \quad f(y) = x, \quad \text{and } f(z) = z \;\; \forall z \in X \setminus \{x,y\}.
	\]
	
	\medskip
	
	\noindent
	\textbf{Example.}  
	Let \(X = \{1,2,3\}\). Define two transpositions:
	\[
	f = (1 \; 2), \qquad g = (2 \; 3).
	\]
	That is,
	\[
	f(1)=2, \; f(2)=1, \; f(3)=3, \qquad 
	g(2)=3, \; g(3)=2, \; g(1)=1.
	\]
	
	Their composition is
	\[
	(f \circ g)(1) = f(g(1)) = f(1) = 2, \quad
	(f \circ g)(2) = f(g(2)) = f(3) = 3, \quad
	(f \circ g)(3) = f(g(3)) = f(2) = 1.
	\]
	
	Thus
	\[
	f \circ g = (1 \; 2 \; 3),
	\]
	which is no longer a transposition, but a $3$-cycle.
	
	
	\subsection*{5. Example Composition}
	
	Take 
	\[
	f = (1\,2), \qquad g = (1\,3\,2).
	\]
	Then
	\[
	\begin{aligned}
		(g \cdot f)(1) &= g(f(1)) = g(2) = 1, \\
		(g \cdot f)(2) &= g(f(2)) = g(1) = 3, \\
		(g \cdot f)(3) &= g(f(3)) = g(3) = 2,
	\end{aligned}
	\]
	so
	\[
	g \cdot f = (2\,3).
	\]
	
	\subsection*{6. Cayley Table of $S_3$}
	
	For convenience, label the elements:
	\[
	e = (), \quad a = (1\,2), \quad b = (2\,3), \quad c = (1\,3), \quad d = (1\,2\,3), \quad d^{-1} = (1\,3\,2).
	\]
	
	The Cayley table under composition is:
	
	\[
	\begin{array}{c|c|c|c|c|c|c}
		\cdot & e & a & b & c & d & d^{-1} \\ \hline
		e & e & a & b & c & d & d^{-1} \\ \hline
		a & a & e & d & d^{-1} & c & b \\ \hline
		b & b & d^{-1} & e & d & a & c \\ \hline
		c & c & d & d^{-1} & e & b & a \\ \hline
		d & d & b & c & a & d^{-1} & e \\ \hline
		d^{-1} & d^{-1} & c & a & b & e & d \\
	\end{array}
	\]
	
	\subsection*{7. Remarks}
	
	- $S_3$ is the smallest non-abelian group.  
	- It has three transpositions, two $3$-cycles, and the identity.  
	- The group law is fully captured by the Cayley table above.
	
	\subsection*{Permutations as Products of Transpositions}
	
	\noindent\textbf{Theorem.}
	Every permutation \(f \in S_n\) can be expressed as a product of transpositions.
	
	\medskip
	\noindent\textbf{Example with \(X=\{1,2,3\}\).}
	The elements of \(S_3\) are
	\(e,\ (12),\ (13),\ (23),\ (123),\ (132)\).
	Here, \((12),(13),(23)\) are transpositions; \((123)\), \((132)\) are $3$-cycles.
	
	\medskip
	\noindent\textbf{General rule for cycles.}
	For a $k$-cycle \((a_1\,a_2\,\dots\,a_k)\),
	\[
	(a_1\,a_2\,\dots\,a_k)=(a_1\,a_k)(a_1\,a_{k-1})\cdots(a_1\,a_2).
	\]
	
	\medskip
	\noindent\textbf{Applying the rule to \((123)\).}
	\[
	(123)=(13)(12).
	\]
	\textit{Check (compose right-to-left):}
	\[
	\begin{array}{c|cc}
		z & (12)\,z & (13)\big((12)\,z\big)\\\hline
		1 & 2 & 3\\
		2 & 1 & 1\\
		3 & 3 & 2
	\end{array}
	\]
	Thus \(1\mapsto2,\ 2\mapsto3,\ 3\mapsto1\), i.e. \((123)\).
	
	\medskip
	\noindent\textbf{Alternative decomposition.}
	\[
	(123)=(12)(23).
	\]
	\textit{Check (compose right-to-left):}
	\[
	\begin{array}{c|cc}
		z & (23)\,z & (12)\big((23)\,z\big)\\\hline
		1 & 1 & 2\\
		2 & 3 & 3\\
		3 & 2 & 1
	\end{array}
	\]
	Again \(1\mapsto2,\ 2\mapsto3,\ 3\mapsto1\), so \((123)\).
	
	
\end{document}
